\section{Method}
\label{sec:method}

\subsection{Problem Setting and Notation}

A Web-RAG pipeline retrieves a candidate set of documents for a query and re-ranks them before generation. We focus on the reranking stage. Let $q$ be the current query, $d$ a candidate document, and $H = \{(q_i, d_i)\}_{i=1}^{K}$ the user's session history of past query--clicked-document pairs. A frozen scorer produces a base relevance score $s_{\mathrm{base}}(q,d) \in \mathbb{R}$ and a document representation $h(q,d) \in \mathbb{R}^{768}$. Our goal is to produce a personalized score $s(q,d,H)$ that re-ranks the candidates in a user-aware manner, using only a small trainable module on top of the frozen features.

\subsection{Pref-Rerank: Bilinear Score Intervention}
\label{sec:method:intervention}

We define the personalized score as
\begin{equation}
\label{eq:prefrerank}
s(q,d,H) \;=\; s_{\mathrm{base}}(q,d) \;+\; h(q,d)^{\top} A\, u(H),
\end{equation}
where $u(H) \in \mathbb{R}^{16}$ is a user vector derived from session history and $A \in \mathbb{R}^{768 \times 16}$ is the only trainable component. Equation~\eqref{eq:prefrerank} is the \texttt{no\_user\_bias} variant reported in our experiments. We also evaluate a \texttt{full} variant that adds a user-only bias term, $s_{\mathrm{base}} + w^\top u + h^\top A u$, with $w \in \mathbb{R}^{16}$; Sec.~\ref{sec:exp:mechanism} shows this term is redundant.

The bilinear term $h^\top A u$ acts as a user-conditioned feature-wise modulation of the document representation: $A u \in \mathbb{R}^{768}$ is a per-user direction in document space, and the score shifts by the alignment of $h(q,d)$ with that direction. This is the standard user--item coupling of factorization machines and bilinear neural recommendation models~\cite{he2017nfm,xiao2017afm}, transferred to evidence reranking. Because $A$ is the only trainable component, Pref-Rerank adds $768 \times 16 = 12{,}288$ parameters (or $12{,}304$ in the \texttt{full} variant), four orders of magnitude fewer than the frozen backbone.

\subsection{User Vector Constructions}
\label{sec:method:uservec}

We study three constructions of $u(H)$, all mapping to $\mathbb{R}^{16}$:

\paragraph{LLM-extracted typed schema (primary).}
A frozen LLM (GPT-5.4-mini via a schema-constrained prompt) reads the first $K=3$ query--clicked-document pairs of the session and outputs a four-dimension closed-vocabulary JSON schema:
\texttt{source} $\in$ \{official\_docs, community, news, academic\},
\texttt{style} $\in$ \{detailed\_specs, concise\_overview, comparison\},
\texttt{evidence\_type} $\in$ \{data\_tables, case\_studies, tutorials\},
\texttt{recency} $\in$ \{latest, any\}.
Each dimension is encoded as a one-hot vector of length 4 (the maximum vocabulary size) and the four blocks are concatenated to form $u \in \{0,1\}^{16}$. The extraction is cached per session history, so training never re-invokes the LLM.

\paragraph{Behavioral session features (ablation).}
We compute six classical session features --- BM25/Jaccard click overlap between the eval query and history queries; a dwell-time proxy (history length); average SERP rank of past clicks; history length; eval query length; and a reformulation flag --- and zero-pad to 16 dimensions. This baseline isolates whether semantic structure (the typed schema) matters beyond behavioral signals.

\paragraph{Random vectors (control).}
We draw $u$ uniformly at random per session (seeded by session id). This control tests whether the bilinear term's capacity, independent of any meaningful user signal, is sufficient to capture session-level patterns.

\subsection{Frozen Backbone: Query-Flow Cross-Encoder}
\label{sec:method:backbone}

AOL-IA provides query and click metadata but not document text. We therefore represent each document by the bag of queries that led to it being clicked anywhere in the corpus --- the query-flow representation~\cite{craswell2007click}. The base score is
\begin{equation}
s_{\mathrm{base}}(q,d) = \mathrm{BM25}\bigl(q,\; \mathrm{concat}(\text{click-queries of } d)\bigr),
\end{equation}
and the document representation is the mean of BGE-base~\cite{xiao2023bge} embeddings of the document's click-queries,
\begin{equation}
h(q,d) = \frac{1}{|Q_d|}\sum_{q' \in Q_d} \mathrm{BGE}(q'), \qquad Q_d = \{\,q' : d \text{ was clicked for } q'\,\}.
\end{equation}
Both $s_{\mathrm{base}}$ and $h$ are precomputed once for every (query, document) pair and frozen during training. This design keeps the trainable surface to a single matrix $A$ and makes experiments reproducible on a single GPU.

\subsection{Training}

\paragraph{Position-debiased supervision.}
We build pairwise training instances from sessions with at least four query events. The first $K=3$ events form the user history; the last event provides the evaluation query. The positive document is the clicked document of the evaluation query, but only if its SERP rank exceeds 3 --- clicks in the top three positions are heavily confounded by visual prominence~\cite{craswell2008experimental,joachims2007evaluating} and are dropped. The negative document is a URL clicked for the \emph{same} query by a \emph{different} anonymized user, excluding any URL in the current user's history.

\paragraph{User-disjoint split.}
Sessions are split into train/dev/test by anonymized user id (70/15/15), so no user appears in more than one split. This prevents the user vector from memorizing user identity across splits.

\paragraph{Objective and optimization.}
We train $\{A, w\}$ (or just $\{A\}$ in \texttt{no\_user\_bias}) by minimizing the Bradley--Terry pairwise ranking loss
\begin{equation}
\mathcal{L} = -\log\sigma\bigl(s(q,d^{+},H) - s(q,d^{-},H)\bigr),
\end{equation}
with AdamW, learning rate $5 \times 10^{-2}$, batch size 64, for 50 epochs, retaining the checkpoint with lowest dev loss. With 12{,}288 parameters and precomputed features, training completes in under one minute on a single RTX 4090.
